\documentclass[10pt,a4paper]{amsart}
\usepackage{amsmath,amssymb,amsthm,graphicx}
% Retained preamble for the Stacks Project chapter snapshot in this directory.
% The chapter's own preamble.tex is a copy of the shared Stacks Project preamble; the block
% below is sliced MECHANICALLY, line for line, from it. Nothing was added, removed or
% redefined; the wrapper class, the xr-hyper/hyperref wiring and the externaldocument list
% are omitted because the assembler supplies the document class.
%
% lines taken from preamble.tex: 3, 5, 9, 11, 12, 156, 157, 158, ... (51 lines total)
%
% package coverage restored from the chapter's own preamble (2026-10-13): xy, xy-2cell

\usepackage[all]{xy}
\xyoption{2cell}
\UseAllTwocells

\newenvironment{reference}{\comment}{\endcomment}
\newenvironment{slogan}{\comment}{\endcomment}
\newenvironment{history}{\comment}{\endcomment}
\theoremstyle{plain}
\newtheorem{theorem}[subsection]{Theorem}
\newtheorem{proposition}[subsection]{Proposition}
\newtheorem{lemma}[subsection]{Lemma}
\theoremstyle{definition}
\newtheorem{definition}[subsection]{Definition}
\newtheorem{example}[subsection]{Example}
\newtheorem{exercise}[subsection]{Exercise}
\newtheorem{situation}[subsection]{Situation}
\theoremstyle{remark}
\newtheorem{remark}[subsection]{Remark}
\newtheorem{remarks}[subsection]{Remarks}
\numberwithin{equation}{subsection}
\def\lim{\mathop{\mathrm{lim}}\nolimits}
\def\colim{\mathop{\mathrm{colim}}\nolimits}
\def\Spec{\mathop{\mathrm{Spec}}}
\def\Hom{\mathop{\mathrm{Hom}}\nolimits}
\def\Ext{\mathop{\mathrm{Ext}}\nolimits}
\def\SheafHom{\mathop{\mathcal{H}\!\mathit{om}}\nolimits}
\def\SheafExt{\mathop{\mathcal{E}\!\mathit{xt}}\nolimits}
\def\Sch{\mathit{Sch}}
\def\Mor{\mathop{\mathrm{Mor}}\nolimits}
\def\Ob{\mathop{\mathrm{Ob}}\nolimits}
\def\Sh{\mathop{\mathit{Sh}}\nolimits}
\def\NL{\mathop{N\!L}\nolimits}
\def\CH{\mathop{\mathrm{CH}}\nolimits}
\def\proetale{{pro\text{-}\acute{e}tale}}
\def\etale{{\acute{e}tale}}
\def\QCoh{\mathit{QCoh}}
\def\Ker{\mathop{\mathrm{Ker}}}
\def\Im{\mathop{\mathrm{Im}}}
\def\Coker{\mathop{\mathrm{Coker}}}
\def\Coim{\mathop{\mathrm{Coim}}}
\DeclareMathSymbol{\boxtimes}{\mathbin}{AMSa}{"02}
\def\QCohstack{\mathcal{QC}\!\mathit{oh}}
\def\Cohstack{\mathcal{C}\!\mathit{oh}}
\def\Spacesstack{\mathcal{S}\!\mathit{paces}}
\def\Quotfunctor{\mathrm{Quot}}
\def\Hilbfunctor{\mathrm{Hilb}}
\def\Curvesstack{\mathcal{C}\!\mathit{urves}}
\def\Polarizedstack{\mathcal{P}\!\mathit{olarized}}
\def\Complexesstack{\mathcal{C}\!\mathit{omplexes}}
\def\Pic{\mathop{\mathrm{Pic}}\nolimits}
\def\Picardstack{\mathcal{P}\!\mathit{ic}}
\def\Picardfunctor{\mathrm{Pic}}
\def\Deformationcategory{\mathcal{D}\!\mathit{ef}}
\usepackage{cleveref}
\crefname{originalreference}{original reference}{original references}
\Crefname{originalreference}{Original reference}{Original references}
\makeatletter
\newlabel{coherent-lemma-chow-Noetherian}{{0200}{source}{}{sourcepointer-1}{}}
\newlabel{coherent-lemma-chow-Noetherian@cref}{{[originalreference][0][]0200}{[1][0][]source}{}{}{}}
\makeatother
\crefname{originalreference}{original reference}{original references}
\Crefname{originalreference}{Original reference}{Original references}
\makeatletter
\newlabel{properties-lemma-point-and-maximal-points-affine}{{01ZX}{source}{}{sourcepointer-2}{}}
\newlabel{properties-lemma-point-and-maximal-points-affine@cref}{{[originalreference][0][]01ZX}{[1][0][]source}{}{}{}}
\makeatother
\crefname{originallemma}{original lemma}{original lemmas}
\Crefname{originallemma}{Original lemma}{Original lemmas}
\makeatletter
\newlabel{lemma-chow-finite-type}{{0202}{source}{}{sourcepointer-3}{}}
\newlabel{lemma-chow-finite-type@cref}{{[originallemma][0][]0202}{[1][0][]source}{}{}{}}
\makeatother
\crefname{originalproposition}{original proposition}{original propositions}
\Crefname{originalproposition}{Original proposition}{Original propositions}
\makeatletter
\newlabel{proposition-separated-closed-in-finite-presentation}{{01ZJ}{source}{}{sourcepointer-4}{}}
\newlabel{proposition-separated-closed-in-finite-presentation@cref}{{[originalproposition][0][]01ZJ}{[1][0][]source}{}{}{}}
\makeatother
\crefname{originalproposition}{original proposition}{original propositions}
\Crefname{originalproposition}{Original proposition}{Original propositions}
\makeatletter
\newlabel{proposition-approximate}{{01ZA}{source}{}{sourcepointer-5}{}}
\newlabel{proposition-approximate@cref}{{[originalproposition][0][]01ZA}{[1][0][]source}{}{}{}}
\makeatother
\crefname{originallemma}{original lemma}{original lemmas}
\Crefname{originallemma}{Original lemma}{Original lemmas}
\makeatletter
\newlabel{lemma-descend-finite-presentation}{{01ZM}{source}{}{sourcepointer-6}{}}
\newlabel{lemma-descend-finite-presentation@cref}{{[originallemma][0][]01ZM}{[1][0][]source}{}{}{}}
\makeatother
\crefname{originallemma}{original lemma}{original lemmas}
\Crefname{originallemma}{Original lemma}{Original lemmas}
\makeatletter
\newlabel{lemma-descend-separated-finite-presentation}{{01ZQ}{source}{}{sourcepointer-7}{}}
\newlabel{lemma-descend-separated-finite-presentation@cref}{{[originallemma][0][]01ZQ}{[1][0][]source}{}{}{}}
\makeatother
\crefname{originallemma}{original lemma}{original lemmas}
\Crefname{originallemma}{Original lemma}{Original lemmas}
\makeatletter
\newlabel{lemma-scheme-over-limit}{{01YZ}{source}{}{sourcepointer-8}{}}
\newlabel{lemma-scheme-over-limit@cref}{{[originallemma][0][]01YZ}{[1][0][]source}{}{}{}}
\makeatother
\crefname{originallemma}{original lemma}{original lemmas}
\Crefname{originallemma}{Original lemma}{Original lemmas}
\makeatletter
\newlabel{lemma-topology-limit}{{01YY}{source}{}{sourcepointer-9}{}}
\newlabel{lemma-topology-limit@cref}{{[originallemma][0][]01YY}{[1][0][]source}{}{}{}}
\makeatother
\crefname{originalreference}{original reference}{original references}
\Crefname{originalreference}{Original reference}{Original references}
\makeatletter
\newlabel{morphisms-definition-scheme-theoretic-image}{{01R7}{source}{}{sourcepointer-10}{}}
\newlabel{morphisms-definition-scheme-theoretic-image@cref}{{[originalreference][0][]01R7}{[1][0][]source}{}{}{}}
\makeatother
\crefname{originalreference}{original reference}{original references}
\Crefname{originalreference}{Original reference}{Original references}
\makeatletter
\newlabel{morphisms-lemma-quasi-compact-scheme-theoretic-image}{{01R8}{source}{}{sourcepointer-11}{}}
\newlabel{morphisms-lemma-quasi-compact-scheme-theoretic-image@cref}{{[originalreference][0][]01R8}{[1][0][]source}{}{}{}}
\makeatother
\crefname{originalreference}{original reference}{original references}
\Crefname{originalreference}{Original reference}{Original references}
\makeatletter
\newlabel{morphisms-lemma-quasi-compact-dominant}{{01RL}{source}{}{sourcepointer-12}{}}
\newlabel{morphisms-lemma-quasi-compact-dominant@cref}{{[originalreference][0][]01RL}{[1][0][]source}{}{}{}}
\makeatother
\title{Chow's lemma for separated finite-type morphisms with finitely many components (EGA route)}
\author{The Stacks Project Authors}
\date{stacks-project:limits.tex:local-snapshot-20261012; GNU FDL 1.2 or later}
\begin{document}
\maketitle
\noindent\textit{Editorial note.} All steps explicit; the alternative 'standard proof' route is mentioned in the source but not required.\par\medskip
\noindent\textit{Allowed standard prerequisites (editorial).}\begin{itemize}
\item Noetherian Chow lemma. Cohomology of Schemes, Lemma coherent-lemma-chow-Noetherian
\end{itemize}\medskip
\ifdefined\hypertarget\hypertarget{sourcepointer-1}{}\fi
\noindent\textit{Source reference.} Original reference 0200: sources/stacks-chow/tags.txt, source lines 2617--2617. Cross-chapter standard prerequisite cited by the author inside the counted span; it is used as a pointer and its own proof is not part of this item.\par\medskip
\ifdefined\hypertarget\hypertarget{sourcepointer-2}{}\fi
\noindent\textit{Source reference.} Original reference 01ZX: sources/stacks-chow/tags.txt, source lines 2612--2612. Cross-chapter standard prerequisite cited by the author inside the counted span; it is used as a pointer and its own proof is not part of this item.\par\medskip
\ifdefined\hypertarget\hypertarget{sourcepointer-3}{}\fi
\noindent\textit{Source reference.} Original lemma 0202: sources/stacks-chow/tags.txt, source lines 2619--2619. Cross-chapter standard prerequisite cited by the author inside the counted span; it is used as a pointer and its own proof is not part of this item.\par\medskip
\ifdefined\hypertarget\hypertarget{sourcepointer-4}{}\fi
\noindent\textit{Source reference.} Original proposition 01ZJ: sources/stacks-chow/tags.txt, source lines 2596--2596. Cross-chapter standard prerequisite cited by the author inside the counted span; it is used as a pointer and its own proof is not part of this item.\par\medskip
\ifdefined\hypertarget\hypertarget{sourcepointer-5}{}\fi
\noindent\textit{Source reference.} Original proposition 01ZA: sources/stacks-chow/tags.txt, source lines 2585--2585. Cross-chapter standard prerequisite cited by the author inside the counted span; it is used as a pointer and its own proof is not part of this item.\par\medskip
\ifdefined\hypertarget\hypertarget{sourcepointer-6}{}\fi
\noindent\textit{Source reference.} Original lemma 01ZM: sources/stacks-chow/tags.txt, source lines 2599--2599. Cross-chapter standard prerequisite cited by the author inside the counted span; it is used as a pointer and its own proof is not part of this item.\par\medskip
\ifdefined\hypertarget\hypertarget{sourcepointer-7}{}\fi
\noindent\textit{Source reference.} Original lemma 01ZQ: sources/stacks-chow/tags.txt, source lines 2603--2603. Cross-chapter standard prerequisite cited by the author inside the counted span; it is used as a pointer and its own proof is not part of this item.\par\medskip
\ifdefined\hypertarget\hypertarget{sourcepointer-8}{}\fi
\noindent\textit{Source reference.} Original lemma 01YZ: sources/stacks-chow/tags.txt, source lines 2574--2574. Cross-chapter standard prerequisite cited by the author inside the counted span; it is used as a pointer and its own proof is not part of this item.\par\medskip
\ifdefined\hypertarget\hypertarget{sourcepointer-9}{}\fi
\noindent\textit{Source reference.} Original lemma 01YY: sources/stacks-chow/tags.txt, source lines 2573--2573. Cross-chapter standard prerequisite cited by the author inside the counted span; it is used as a pointer and its own proof is not part of this item.\par\medskip
\ifdefined\hypertarget\hypertarget{sourcepointer-10}{}\fi
\noindent\textit{Source reference.} Original reference 01R7: sources/stacks-chow/tags.txt, source lines 2292--2292. Cross-chapter standard prerequisite cited by the author inside the counted span; it is used as a pointer and its own proof is not part of this item.\par\medskip
\ifdefined\hypertarget\hypertarget{sourcepointer-11}{}\fi
\noindent\textit{Source reference.} Original reference 01R8: sources/stacks-chow/tags.txt, source lines 2293--2293. Cross-chapter standard prerequisite cited by the author inside the counted span; it is used as a pointer and its own proof is not part of this item.\par\medskip
\ifdefined\hypertarget\hypertarget{sourcepointer-12}{}\fi
\noindent\textit{Source reference.} Original reference 01RL: sources/stacks-chow/tags.txt, source lines 2306--2306. Cross-chapter standard prerequisite cited by the author inside the counted span; it is used as a pointer and its own proof is not part of this item.\par\medskip
\begin{lemma}
\label{lemma-chow-EGA}
Let $S$ be a quasi-compact and quasi-separated scheme.
Let $f : X \to S$ be a separated morphism of finite type.
Assume that $X$ has finitely many irreducible components.
Then there exists an $n \geq 0$ and a diagram
$$
\xymatrix{
X \ar[rd] & X' \ar[d] \ar[l]^\pi \ar[r] & \mathbf{P}^n_S \ar[dl] \\
& S &
}
$$
where $X' \to \mathbf{P}^n_S$ is an immersion, and
$\pi : X' \to X$ is proper and surjective. Moreover, there exists
an open dense subscheme $U \subset X$ such that $\pi^{-1}(U) \to U$
is an isomorphism of schemes.
\end{lemma}

\begin{proof}
Let $X = Z_1 \cup \ldots \cup Z_n$ be the decomposition of $X$
into irreducible components. Let $\eta_j \in Z_j$ be the generic point.

\medskip\noindent
There are (at least) two ways to proceed with the proof.
The first is to redo the proof of
Cohomology of Schemes, Lemma \ref{coherent-lemma-chow-Noetherian}
using the general
Properties, Lemma \ref{properties-lemma-point-and-maximal-points-affine}
to find suitable affine opens in $X$. (This is the ``standard'' proof.)
The second is to use absolute Noetherian approximation as in
the proof of Lemma \ref{lemma-chow-finite-type} above.
This is what we will do here.

\medskip\noindent
By Proposition \ref{proposition-separated-closed-in-finite-presentation}
we can find a closed immersion $X \to Y$ where $Y$ is separated
and of finite presentation over $S$.
Write $S = \lim_i S_i$ as a directed limit of Noetherian schemes, see
Proposition \ref{proposition-approximate}. By
Lemma \ref{lemma-descend-finite-presentation} we can
find an index $i \in I$ and a scheme $Y_i \to S_i$ of finite presentation
so that $Y = S \times_{S_i} Y_i$.
By Lemma \ref{lemma-descend-separated-finite-presentation}
we may assume that $Y_i \to S_i$ is separated.
We have the following diagram
$$
\xymatrix{
\eta_j \in Z_j \ar[r] & X \ar[r] \ar[rd] & Y \ar[r] \ar[d] & Y_i \ar[d] \\
& & S \ar[r] & S_i
}
$$
Denote $h : X \to Y_i$ the composition.

\medskip\noindent
For $i' \geq i$ write $Y_{i'} = S_{i'} \times_{S_i} Y_i$.
Then $Y = \lim_{i' \geq i} Y_{i'}$, see
Lemma \ref{lemma-scheme-over-limit}.
Choose $j, j' \in \{1, \ldots, n\}$, $j \not = j'$.
Note that $\eta_j$ is not a specialization of $\eta_{j'}$.
By Lemma \ref{lemma-topology-limit}
we can replace $i$ by a bigger index and assume
that $h(\eta_j)$ is not a specialization of $h(\eta_{j'})$
for all pairs $(j, j')$ as above.
For such an index, let
$Y' \subset Y_i$ be the scheme theoretic image of
$h : X \to Y_i$, see
Morphisms, Definition \ref{morphisms-definition-scheme-theoretic-image}.
The morphism $h$ is quasi-compact as the composition of the quasi-compact
morphisms $X \to Y$ and $Y \to Y_i$ (which is affine).
Hence by
Morphisms, Lemma \ref{morphisms-lemma-quasi-compact-scheme-theoretic-image}
the morphism $X \to Y'$ is dominant. Thus the generic points
of $Y'$ are all contained in the set
$\{h(\eta_1), \ldots, h(\eta_n)\}$, see
Morphisms, Lemma \ref{morphisms-lemma-quasi-compact-dominant}.
Since none of the $h(\eta_j)$ is the specialization of another
we see that the points $h(\eta_1), \ldots, h(\eta_n)$ are pairwise
distinct and are each a generic point of $Y'$.

\medskip\noindent
We apply Cohomology of Schemes, Lemma
\ref{coherent-lemma-chow-Noetherian} above to the morphism
$Y' \to S_i$. This gives a diagram
$$
\xymatrix{
Y' \ar[rd] & Y^* \ar[d] \ar[l]^\pi \ar[r] & \mathbf{P}^n_{S_i} \ar[dl] \\
& S_i &
}
$$
such that $\pi$ is proper and surjective and an isomorphism over
a dense open subscheme $V \subset Y'$. By our choice of $i$ above
we know that $h(\eta_1), \ldots, h(\eta_n) \in V$. Consider
the commutative diagram
$$
\xymatrix{
X' \ar@{=}[r] &
X \times_{Y'} Y^* \ar[r] \ar[d] &
Y^* \ar[r] \ar[d] &
\mathbf{P}^n_{S_i} \ar[ddl] \\
& X \ar[r] \ar[d] & Y' \ar[d] & \\
& S \ar[r] & S_i &
}
$$
Note that $X' \to X$ is an isomorphism over the open subscheme
$U = h^{-1}(V)$ which contains each of the $\eta_j$ and hence is
dense in $X$. We conclude $X \leftarrow X' \rightarrow \mathbf{P}^n_S$
is a solution to the problem posed in the lemma.
\end{proof}

\end{document}
